\section{Hướng phát triển}

Hướng phát triển được xây dựng từ nền tảng hiện tại như
Hình~\ref{fig:future-directions}.

\begin{figure}[H]
\centering
\resizebox{0.92\textwidth}{!}{%
\begin{tikzpicture}[node distance=0.65cm and 0.8cm]
  \node[groupbox, text width=5cm] (current) {Hiện tại\\YOLO Pruning + QAT\\Jetson\\Edge Platform\\Copilot};
  \node[groupbox, right=2cm of current, text width=6cm] (future) {Tương lai\\Federated Learning\\Multi-node Orchestration\\Advanced Anomaly Detection\\Auto-remediation\\Model A/B Testing\\More Edge Hardware\\Autonomous Edge AI Operations};
  \draw[flowarrow] (current) -- (future);
\end{tikzpicture}}
\caption{Các hướng phát triển từ kết quả hiện tại}
\label{fig:future-directions}
\end{figure}

Thứ nhất, nền tảng có thể mở rộng sang nhiều device class và tự động chọn model
artifact theo capability, power budget và mục tiêu latency. Model A/B testing và
shadow deployment giúp so sánh phiên bản mới trên dữ liệu thực trước khi chuyển
toàn bộ lưu lượng.

Thứ hai, multi-node orchestration, quota và scheduler có thể tối ưu vị trí chạy
runtime trên một cụm Edge. Federated learning hoặc continuous learning cho phép
cải thiện model từ dữ liệu phân tán nhưng phải bổ sung quản trị dữ liệu, bảo vệ
riêng tư và kiểm định model trước phát hành.

Thứ ba, Monitoring có thể phát triển từ rule-based sang anomaly detection đa biến,
drift detection và dự báo tài nguyên. Incident kết hợp runbook machine-readable
có thể hỗ trợ auto-remediation theo phạm vi giới hạn, rollback tự động và kiểm tra
hậu điều kiện.

Cuối cùng, Copilot có thể tiến tới Autonomous Edge AI Operations theo từng cấp
độ tự chủ. Mọi mức tự động hóa cao hơn phải đi kèm sandbox, policy, simulation,
approval threshold, audit và cơ chế dừng khẩn cấp để bảo đảm con người vẫn kiểm
soát các hành động rủi ro.